\documentclass[11pt]{article}

\usepackage[margin=1in]{geometry}
\usepackage[T1]{fontenc}
\usepackage{lmodern}
\usepackage{microtype}
\usepackage{array}
\usepackage{longtable}
\usepackage{booktabs}
\usepackage{amsmath}
\usepackage{amsfonts}
\usepackage{xcolor}
\usepackage{url}

\definecolor{donecolor}{RGB}{20,105,62}
\definecolor{opencolor}{RGB}{160,55,35}
\definecolor{conditionalcolor}{RGB}{135,92,18}
\newcommand{\lean}[1]{\path{#1}}
\newcommand{\done}{\textcolor{donecolor}{locally complete}}
\newcommand{\openstate}[1]{\textcolor{opencolor}{open (#1)}}
\newcommand{\conditional}{\textcolor{conditionalcolor}{conditional}}

\makeatletter
\renewcommand\section{\@startsection{section}{1}{0pt}%
  {2.4ex plus .6ex minus .2ex}{1.0ex}%
  {\normalfont\Large}}
\renewcommand\subsection{\@startsection{subsection}{2}{0pt}%
  {1.8ex plus .4ex minus .2ex}{0.7ex}%
  {\normalfont\large}}
\makeatother

\setlength{\parindent}{0pt}
\setlength{\parskip}{0.55em}
\renewcommand{\arraystretch}{1.16}
\renewcommand{\descriptionlabel}[1]{\hspace{\labelsep}\normalfont #1}

\begin{document}

\begin{center}
  {\LARGE State and proof roadmap for \lean{Lemmas}}\\[0.6em]
  {\large Quantitative strict-AT formalization}\\[0.4em]
  Snapshot checked on 1 August 2026
\end{center}

\section{Executive status}

The library target \lean{LatalaMeetsAT} builds successfully with Lean 4.32.1 and
mathlib 4.32.1.  The directory \lean{Lemmas} contains 26 Lean files.  Fourteen
files have no local \lean{sorry}; twelve files contain 18 occurrences in total.
There are no project-specific \lean{axiom} declarations and no \lean{admit}s in
this directory.

The build is therefore structurally complete but not certified.  Lean reports
that \lean{SpinGlass.AT.quantitative_strictAT} depends on
\lean{sorryAx}, in addition to the usual foundations
\lean{propext}, \lean{Classical.choice}, and \lean{Quot.sound}.  A file marked
``locally complete'' below has no \lean{sorry} in that file; it may still depend
transitively on an open declaration imported earlier.

Four of the 18 holes are definitions rather than propositions:
\lean{replicaCovarianceOperator}, \lean{scalarOrderParameter}, and
\lean{gtSemigroupSolution}, together with
\lean{cavityInterpolatedAverage}.  These are foundational interfaces.  Filling
later theorem holes before fixing these interfaces would hide the main analytic
work rather than remove it.

\section{Import spine}

The current imports form one linear spine.  This determines build order, while
the mathematical work can still be split into coherent packages.

\begin{center}
\begin{minipage}{0.94\textwidth}\small
\lean{Replicas}
$\to$ \lean{RSParameters}
$\to$ \lean{UniformData}
$\to$ \lean{SmartPath}
$\to$ \lean{GaussianDerivative}
$\to$ \lean{FreeEnergyDerivative}
$\to$ \lean{Scalar/Semigroup}
$\to$ \lean{Scalar/LatalaKernel}
$\to$ \lean{Scalar/StrictATSign}
$\to$ \lean{GT/Defs}
$\to$ \lean{GT/Semigroup}
$\to$ \lean{GT/HalfCascade}
$\to$ \lean{GT/Interpolation}
$\to$ \lean{GT/Coercivity}
$\to$ \lean{CoupledPressure}
$\to$ \lean{FixedDeviation}
$\to$ \lean{Cavity/Defs}
$\to$ \lean{Cavity/Interpolation}
$\to$ \lean{Cavity/Coefficients}
$\to$ \lean{Cavity/System}
$\to$ \lean{Cavity/Stability}
$\to$ \lean{Absorption}
$\to$ \lean{FreeEnergy}
$\to$ \lean{Replicon}
$\to$ \lean{CompactCorollary}
$\to$ \lean{MainResult}.
\end{minipage}
\end{center}

\section{State of every Lean file}

\small
\begin{longtable}{@{}p{0.28\textwidth}p{0.15\textwidth}p{0.50\textwidth}@{}}
\toprule
File & Local state & What is present or missing \\
\midrule
\endhead
\lean{Replicas.lean} & \done & Finite configurations, Gibbs and quenched replica averages, overlap bounds, linearity, monotonicity, and the mixed-overlap estimate. \\
\lean{RSParameters.lean} & \done & Gaussian RS quantities, the selected fixed point, positivity and range of $q$, identities for $A$, and anomalous-versus-replicon comparison. \\
\lean{UniformData.lean} & \done & Compact-uniform parameter package and the strict-AT path gap. \\
\lean{SmartPath.lean} & \done & Covariance kernel, abstract centered Gaussian path, full Hamiltonian with external field, and path gap.  An explicit finite Gaussian realization is not supplied. \\
\lean{GaussianDerivative.lean} & \openstate{2} & Covariance operator and the reusable Gaussian differentiation theorem remain. \\
\lean{FreeEnergyDerivative.lean} & \openstate{1} & Definitions and \lean{rsGap_deriv} are complete; \lean{smartPath_freeEnergy_deriv} remains. \\
\lean{Scalar/Semigroup.lean} & \openstate{1} & The two branches of \lean{scalarPsi} are proved; identification of the scalar trial value with the RS path remains. \\
\lean{Scalar/LatalaKernel.lean} & \openstate{3} & The derivative and sign of \lean{latalaH} and the abstract opposite-monotonicity covariance inequality are complete.  Kernel antitonicity, normalization, and the weighted comparison remain. \\
\lean{Scalar/StrictATSign.lean} & \openstate{2} & The scalar diffusion order parameter and all strict sign estimates remain. \\
\lean{GT/Defs.lean} & \done & Attainable overlaps, constrained free energy, signed matrix path, mass profile, correction, terminal condition, and endpoint identities. \\
\lean{GT/Semigroup.lean} & \done & Elementary zero, half, and one semigroup definitions; only continuity of the zero operator is proved. \\
\lean{GT/HalfCascade.lean} & \done & Constant-function half-cascade identity. \\
\lean{GT/Interpolation.lean} & \openstate{2} & The two-dimensional recursion and the two-replica Guerra--Talagrand bound remain. \\
\lean{GT/Coercivity.lean} & \openstate{1} & The trivial positive local constant is proved; uniform quadratic coercivity remains. \\
\lean{CoupledPressure.lean} & \openstate{1} & Correct quadratic pressure definitions are present.  The uniform sublinear overlap estimate remains; its public alias is complete. \\
\lean{FixedDeviation.lean} & \openstate{1} & Exponentially small fixed-deviation tails remain. \\
\lean{Cavity/Defs.lean} & \done & The observables $A,B,C$, cavity vector, source vector, and $3\times3$ cavity matrix. \\
\lean{Cavity/Interpolation.lean} & \openstate{2} & The elementary third-moment bound is complete.  The last-spin interpolation and its second-derivative estimate remain. \\
\lean{Cavity/Coefficients.lean} & \done & Replica-edge classification, decoupled coefficients, and equality with the displayed cavity matrix. \\
\lean{Cavity/System.lean} & \conditional & No local hole.  The system is an algebraic identity because the remainder is defined as the residual.  Its analytic bound is a predicate and is assumed, not constructed. \\
\lean{Cavity/Stability.lean} & \done & Determinant formula, anomalous gap, uniform inverse estimate, and replicon left eigenvector. \\
\lean{Absorption.lean} & \openstate{1} & Norm domination is complete; the uniform $N A_s$ bound remains. \\
\lean{FreeEnergy.lean} & \openstate{1} & RS and finite-volume free energies are defined; the uniform $O(N^{-1})$ error remains. \\
\lean{Replicon.lean} & \done & Third-moment little-$o$ and replicon susceptibility are proved from the earlier tail, absorption, stability, and assumed cavity-remainder results. \\
\lean{CompactCorollary.lean} & \done & Compactness plus continuity and strict AT produce \lean{UniformATData}. \\
\lean{MainResult.lean} & \conditional & The three conclusions are assembled correctly, conditional on an existential \lean{HasCavityRemainderBound}; transitive holes yield \lean{sorryAx}. \\
\bottomrule
\end{longtable}
\normalsize

\section{Difficulty and recommended prover for open files}

The rating is for the file in its present repository state, including missing
helper lemmas and definitions.  ``Codex'' means it should lead the work when the
task requires designing definitions, repairing APIs, adding supporting lemmas,
or coordinating changes across files.  ``Aristotle'' means the statement and
interfaces are already stable enough for focused Lean proof search.  A handoff
to Aristotle is useful only after the named prerequisite has been supplied.
No open file is rated easy at present.

\small
\begin{longtable}{@{}p{0.25\textwidth}p{0.12\textwidth}p{0.13\textwidth}p{0.42\textwidth}@{}}
\toprule
File & Difficulty & Better lead & Reason and possible handoff \\
\midrule
\endhead
\lean{GaussianDerivative.lean} & difficult & Codex & A core definition and a finite-Gaussian differentiation API are absent.  Hand off only the final finite-sum and integration-by-parts proof after that API is fixed. \\
\lean{FreeEnergyDerivative.lean} & difficult & Codex & The proof must instantiate the Gaussian API, derive the covariance cancellation, and normalize replica sums.  Aristotle can help with the final algebra after those identities are named. \\
\lean{Scalar/Semigroup.lean} & moderate & Codex & The main missing tool is a properly stated Gaussian convolution lemma with integrability conditions.  Once added, Aristotle is well suited to the remaining rewriting and ring normalization. \\
\lean{Scalar/LatalaKernel.lean} & difficult & Codex & The file combines parameter differentiation under an expectation, an endpoint-singular reference density, and a normalization change of variables.  Aristotle is a good follow-up for \lean{latala_weighted_kernel_le} after the first two holes close. \\
\lean{Scalar/StrictATSign.lean} & difficult & Codex & The order parameter itself is missing and requires a PDE or stochastic-process interface.  This is definition and architecture work before proof search can begin. \\
\lean{GT/Interpolation.lean} & difficult & Codex & The two-dimensional recursion, cascade construction, endpoint equations, and derivative sign must be built together.  Aristotle can close local endpoint or algebraic sublemmas once extracted. \\
\lean{GT/Coercivity.lean} & difficult & Codex & The proof needs new multiplier-derivative, Taylor, compactness, and signed-range lemmas.  Its present single theorem is too broad for isolated proof search. \\
\lean{CoupledPressure.lean} & difficult & Codex & The argument coordinates GT coercivity, overlap discretization, Gaussian concentration, convexity, Gronwall, and a uniform limit sequence. \\
\lean{FixedDeviation.lean} & difficult & Codex & The abstract disorder lacks the explicit coordinate or concentration interface needed for the gradient estimate.  The model-level gap must be repaired first. \\
\lean{Cavity/Interpolation.lean} & difficult & Codex & The last-spin interpolation is undefined, and the derivative expansion requires systematic replica bookkeeping and uniform estimates.  Aristotle can handle bounded coefficient identities after they are separated. \\
\lean{Absorption.lean} & moderate & Aristotle & Conditional on the completed tail and cavity-remainder estimates, the target and required bounds are already clear.  The remaining work is chiefly inequalities, constant choices, and finite-size cleanup. \\
\lean{FreeEnergy.lean} & moderate & Codex & The source still needs endpoint and closed-interval calculus lemmas.  After Codex adds those interfaces, Aristotle can finish nonnegativity and the uniform $M/N$ estimate. \\
\bottomrule
\end{longtable}
\normalsize

\section{The 18 remaining sorry obligations}

\subsection{Gaussian interpolation and free-energy derivative}

\begin{enumerate}
\item \lean{GaussianDerivative.lean:replicaCovarianceOperator}.  Give the
explicit finite-replica covariance-derivative expression.  It should be stated
independently of \lean{deriv}, so that the next theorem is not circular.

\item \lean{GaussianDerivative.lean:quenchedGibbs_deriv_of_covariance_deriv}.
Prove differentiability of the quenched finite Gibbs expectation along the
Gaussian covariance path and identify its derivative with the operator above.
The missing reusable ingredient is uniqueness and smooth dependence of a
finite-dimensional centered Gaussian law on its covariance, or an explicit
affine Gaussian realization with the required domination estimates.

\item \lean{FreeEnergyDerivative.lean:smartPath_freeEnergy_deriv}.  Differentiate
the finite log partition function, apply Gaussian integration by parts, prove
the finite-volume covariance cancellation, and rewrite the replica sum as
\lean{overlapSecondMoment}.  This should give
$\frac{\beta^2}{4}((1-q)^2-\mathbb E\langle(R_{12}-q)^2\rangle)$.
The already proved \lean{rsGap_deriv} then follows without a new analytic
argument.
\end{enumerate}

\subsection{Scalar semigroup and Lata\l a comparison}

\begin{enumerate}
\setcounter{enumi}{3}
\item \lean{Scalar/Semigroup.lean:scalarTrialValue_eq}.  Evaluate the lower
branch at $u=0$, combine the two independent Gaussian contributions, and
simplify the deterministic correction to \lean{rsPathValue}.  The proof needs
the Gaussian convolution identity and its integrability hypotheses.

\item In \lean{Scalar/LatalaKernel.lean}, prove \lean{latalaF_antitone}.  Lift the pointwise
nonpositivity of $\partial_y H(t,y)$ through the Gaussian expectation.  Establish
measurability, integrability, and positivity of the factors on $[0,1]$.

\item In \lean{Scalar/LatalaKernel.lean}, prove the reference-mean identity
\lean{latalaF_reference_mean_one}:
$\int_0^1 F_\lambda(y)/(2\sqrt{1-y})\,dy=1$.  This is the central normalization
identity and needs the paper's change of variables or an equivalent Gaussian
integral calculation, including the integrable endpoint singularity at $y=1$.

\item \lean{Scalar/LatalaKernel.lean:latala_weighted_kernel_le}.  Convert the
weighted interval integrals to a probability measure with density
\lean{referenceDensity}, apply the completed opposite-monotonicity covariance
inequality, then use the two normalization assumptions.  Any missing
measurability and restricted-to-global monotonicity lemmas should be stated
explicitly.

\item \lean{Scalar/StrictATSign.lean:scalarOrderParameter}.  Construct
$g_s(u)=\mathbb E[(\partial_x\Psi_s(u,X_u))^2]$ for the local-field diffusion.
This requires the scalar PDE or stochastic-process interface, not a placeholder
real number.

\item \lean{Scalar/StrictATSign.lean:strictAT_sign}.  Below $q$, use It\^o and
conditional Jensen to compare $g_s'$ with the path AT parameter.  Above $q$,
use the Lata\l a kernel estimate.  Integrate the slope inequalities and use
compactness to obtain the uniform local constants $c$ and $\varepsilon$.
\end{enumerate}

\subsection{Guerra--Talagrand interpolation}

\begin{enumerate}
\setcounter{enumi}{9}
\item \lean{GT/Interpolation.lean:gtSemigroupSolution}.  Define the finite
two-dimensional Parisi recursion for \lean{signedMatrixPath},
\lean{gtMassParameter}, and \lean{gtTerminal}; prove the domain, integrability,
and endpoint properties used later.

\item \lean{GT/Interpolation.lean:twoReplica_GT_bound}.  Build the finite
half-mass cascade, identify both endpoints, and prove the Gaussian interpolation
derivative is nonpositive.  At multiplier zero, prove factorization into the
two scalar RS semigroups and match \lean{gtCorrection}.

\item \lean{GT/Coercivity.lean:gt_quadratic_coercivity}.  Differentiate the GT
functional in its multiplier.  Use \lean{strictAT_sign} and a uniform second
derivative bound for the local quadratic gap; use compactness for overlaps away
from $q$ and handle both signed-overlap ranges.  Convert the fixed negative gap
to $-c(v-q)^2$ using $|v-q|\leq2$.
\end{enumerate}

\subsection{Coupled pressure and concentration}

\begin{enumerate}
\setcounter{enumi}{12}
\item \lean{CoupledPressure.lean:coupledPressure_sublinear}.  Choose a coupling
smaller than the GT coercivity constant, split over the $N+1$ attainable
overlaps, and control the maximum centered constrained pressure by Gaussian
concentration.  Combine this envelope with \lean{rsGap_deriv}, convexity in the
coupling, and Gronwall.  Construct an explicit sequence such as
$C\sqrt{\log(N+1)/N}+C\log(N+1)/N$ and prove that it tends to zero uniformly.

\item \lean{FixedDeviation.lean:fixedDeviation}.  For the quadratic exponential
moment, prove the $O_K(N)$ squared-gradient bound and apply Gaussian
concentration.  Combine it with the pointwise exponential Markov estimate and
absorb finitely many small sizes into the prefactor.  This proof needs an
explicit finite Gaussian coordinate model, or a concentration theorem stated
directly for the abstract Gaussian vector in \lean{RSSmartPathDisorder}.
\end{enumerate}

\subsection{Cavity interpolation, absorption, and free energy}

\begin{enumerate}
\setcounter{enumi}{14}
\item \lean{Cavity/Interpolation.lean:cavityInterpolatedAverage}.  Define the
last-spin interpolation $\nu_{s,u}(F)$ with the correct coupled disorder,
normalization, and replica count.  State measurability and endpoint identities
needed by the cavity expansion.

\item \lean{Cavity/Interpolation.lean:cavity_secondDerivative_bound}.  Differentiate
the last-spin interpolation twice, classify replica-edge terms using
\lean{Cavity/Coefficients}, and bound every remainder uniformly by
$C_K(\operatorname{thirdMoment}+N^{-3})$.

\item \lean{Absorption.lean:uniform_secondMoment}.  Apply the cavity system and
the uniform inverse bound.  Split $\mathbb E|R_{12}-q|^3$ at a small threshold,
use \lean{fixedDeviation} for the large-deviation part, and absorb the small
multiple of $A_s$ to obtain $N A_s\leq M$ uniformly.

\item \lean{FreeEnergy.lean:rs_freeEnergy_error}.  Prove the product endpoint
identity at $s=0$, supply continuity and integrability on the closed interval,
and integrate \lean{smartPath_freeEnergy_deriv}.  Nonnegativity of the overlap
second moment gives the lower bound, while \lean{uniform_secondMoment} and the
uniform $\beta$ bound give the $M/N$ upper bound.
\end{enumerate}

\section{Recommended proof roadmap}

\begin{description}
\item[Model and Gaussian API.] Add a finite-index Gaussian-law module.  Decide
whether the current smart-path structure remains law-based or gains an explicit affine
realization.  Finish obligations 1--3 and add endpoint lemmas for
\lean{fullPathHamiltonian}.  This removes the common blocker for free-energy,
pressure, and concentration proofs.

\item[Scalar package.] Finish obligations 4--9.  First isolate Gaussian
convolution and the reference-measure normalization as reusable lemmas.  Then
construct the PDE/diffusion order parameter and prove the uniform strict sign.

\item[GT package.] Finish obligations 10--12.  The recursion should expose its
multiplier derivatives and factorization theorem as named lemmas, since
coercivity uses both.

\item[Pressure and tails.] Finish obligations 13--14.  Reuse the same finite
Gaussian concentration interface created for the smart path, rather than
building a second probabilistic model.

\item[Cavity package.] Finish obligations 15--16 and connect the Taylor
remainder to \lean{cavityRemainder}.  Prove a theorem that constructs
\lean{HasCavityRemainderBound} from the interpolation estimate.  This connection
is currently absent and is necessary to remove the hypothesis from a fully
unconditional version of the main theorem.

\item[Absorption and integration.] Finish obligations 17--18.  The existing
\lean{Replicon}, \lean{CompactCorollary}, and \lean{MainResult} proofs should
then close by recompilation, apart from the explicit cavity-remainder hypothesis
and any intended SK bridge assumptions outside \lean{Lemmas}.

\item[Certification.] Run \lean{lake build LatalaMeetsAT}, search the source for
\lean{sorry}, \lean{admit}, and project-specific \lean{axiom}, and inspect
\lean{#print axioms SpinGlass.AT.quantitative_strictAT}.  Success means that
\lean{sorryAx} is absent.  Ordinary foundations such as choice and quotient
soundness may remain.
\end{description}

\section{Structural issues not counted among the 18 holes}

\begin{itemize}
\item \lean{Cavity/System.lean:cavity_system} is true by unfolding the
definition of \lean{cavityRemainder}.  It does not itself formalize the cavity
expansion.  The analytic content is precisely the unconstructed bound described
above.

\item \lean{MainResult.lean:quantitative_strictAT} assumes the existence of a
constant satisfying \lean{HasCavityRemainderBound data C}.  Thus removing all
18 occurrences is enough to remove the present explicit placeholders, but not
enough to derive that hypothesis from the compact strict-AT assumptions.

\item The free-energy comment saying that the external field is omitted is
stale.  The definition \lean{pathFreeEnergy} now uses
\lean{fullPathHamiltonian}, which adds the deterministic field.  What is still
missing is a proved endpoint identity and the closed-interval calculus needed
to integrate the derivative.

\item \lean{SmartPath.lean} specifies Gaussian marginals and their covariance,
but not a concrete coordinate realization.  This is adequate only after proving
a finite-dimensional Gaussian-law uniqueness and concentration API.  Otherwise
the structure should be strengthened once, near the root of the import spine.
\end{itemize}

\end{document}
